Just like the previous versions, v1 generated lots of learning outcomes and progressed us further towards a fully-functional camera payload. The most notable change to this version was the top-down systems-engineering approach, defining system requirements and completing trade studies to ensure the selected hardware was appropriate. This method heavily reduced the additional scope to the version and allowed me to focus on the core parts of the project without running off on tangents. The board was designed with HAB and CubeSat integration in-mind, which also helped define requirements and PCB design. Mistakes from v0 and v0.1 also helped me develop the hardware efficiently and in an organised matter. Moreover, testing the camera module on a breadboard before procuring the PCB was a good decision as I could verify whether my pinout was correct. I was very happy with the hardware aspect of the project, as the mistakes I made from v0 and v0.1 made it clear planning could not be skipped. Unfortunately, I did not execute the same level of consideration with the sofware which generated most of the issues with this version. This oversight could be explained by the fact I didn't fully explore the software/programming of the previous versions so there weren't as many mistakes to learn from, regardless, these serve as lessons for future versions.
\nl
As I developed the hardware, my design choices that were justified by software assumed the implementation could be executed e.g., I can run the CSA and camera at the same time because I should be able to implement an RTOS. These assumptions eventually led to oversights like not being able to transfer the DCMI image directly into the external memory as I assumed the GPDMA could do it. Although these mistakes were not fatal to the version, it highlights the importance of planning in all aspects of the development cycle, not just the hardware, meaning hardware choices should be backed up by system requirements and software capabilities before being executed. To summarise the assumptions I had made:

\begin{itemize}
    \item The MCU could be placed in to standby mode and woken up from any external interrupt pin.
    \item I could write to the microSD card using transmit only SPI.
    \item The GPDMA could move data from the DCMI peripheral to the QSPI peripheral directly.
    \item That ST natively supported USB FS (but "USB Device" lost support on the U5).
    \item I could transmit image data with errors using USB Communications Device Class.
\end{itemize}

Despite the errors, the version still proved that it is possible to capture and transmit image data to a dashboard, effectively succeeding as a prototype for the CubeSat payload. Future versions of this board should implement the above suggestions and integrate CubeSat and HAB requirements into the system requirements, like those from the Phase A and B work packages, as well as testing the board with other subsystems like with a mechanical chassis and EPS.